\section{Audit design}\label{sec:method}

The audit combines four forms of evidence (Fig.~\ref{fig:overview}): a
mechanistic analysis of the deployed score, external validation against human
opinion scores, construction-defined degradation pairs in the broadcast domain,
and an operation-level substitution in the summarization pipeline. A method is
eligible only if it responds to controlled degradation, generalizes to every
unseen opinion database and exceeds a reference-model agreement floor in the
deployment domain.

\begin{figure}[!ht]
\centering
\resizebox{\textwidth}{!}{%
\begin{tikzpicture}[
  stage/.style={rectangle,rounded corners=2.2pt,minimum height=.68cm,
    text width=2.05cm,align=center,font=\sffamily\scriptsize,
    draw=rulegray!85,fill=gray!5,line width=.55pt,inner xsep=3pt},
  focus/.style={stage,draw=catbovw,fill=catbovw!10,line width=1.15pt,
    font=\sffamily\bfseries\scriptsize},
  operation/.style={rectangle,rounded corners=3pt,minimum height=.78cm,
    text width=11.55cm,align=center,font=\sffamily\scriptsize,
    draw=catbovw,fill=catbovw!6,line width=.9pt,inner ysep=4pt},
  audit/.style={rectangle,rounded corners=2.5pt,minimum height=1.05cm,
    text width=5.25cm,align=left,font=\sffamily\scriptsize,
    draw=rulegray!90,fill=gray!3,line width=.65pt,inner xsep=7pt,inner ysep=5pt},
  evidencehead/.style={rectangle,rounded corners=5pt,minimum height=.38cm,
    font=\sffamily\bfseries\tiny,text=white,fill=headerblue,draw=headerblue,
    inner xsep=9pt,inner ysep=2.5pt},
  flow/.style={-{Stealth[length=4.8pt,width=4.2pt]},line width=.65pt,
    draw=rulegray!95,shorten <=1pt,shorten >=1pt},
  focusflow/.style={-{Stealth[length=5.2pt,width=4.6pt]},line width=.9pt,
    draw=catbovw,shorten <=1pt,shorten >=1pt}]

\node[font=\sffamily\bfseries\tiny,text=headerblue,anchor=west]
  at (-6.15,5.05) {HSMVIDEOSUMM PIPELINE};
\node[stage] (a) at (-5.10,4.47) {Introductory\\passages};
\node[stage] (b) at (-2.55,4.47) {Subjectivity\\filtering};
\node[stage] (c) at (0,4.47) {Redundancy\\detection};
\node[focus] (d) at (2.55,4.47) {Visual-quality\\selection};
\node[stage] (e) at (5.10,4.47) {Chronological\\ordering};
\draw[flow] (a.east)--(b.west);
\draw[flow] (b.east)--(c.west);
\draw[flow] (c.east)--(d.west);
\draw[flow] (d.east)--(e.west);

\node[operation] (op) at (0,3.20)
  {\textcolor{catbovw}{\textbf{AUDITED STAGE}}\enspace Visual-quality selection\\[1pt]
   \textbf{Criterion:} visual quality\qquad
   \textbf{Operation:} SIFT $\rightarrow$ BoVW $\rightarrow$ TF-IDF sum};
\draw[focusflow] (d.south)--(d.south |- op.north);

\node[evidencehead] (eh) at (0,2.30) {FOUR COMPLEMENTARY TESTS};
\draw[focusflow] (op.south)--(eh.north);

\node[audit] (m) at (-3.0,1.43)
  {\textcolor{headerblue}{\textbf{01\enspace Mechanism}}\\[2pt]
   Two IDF forms; count $\times$ mean IDF\\weight; 19 extractor substitutions.};
\node[audit] (x) at (3.0,1.43)
  {\textcolor{headerblue}{\textbf{02\enspace External validation}}\\[2pt]
   KADID, KonIQ, CLIVE; MOS\\and controlled degradation levels.};
\node[audit] (g) at (-3.0,.10)
  {\textcolor{headerblue}{\textbf{03\enspace Constructed truth}}\\[2pt]
   20 broadcast segments; 800\\original--degraded paired decisions.};
\node[audit] (p) at (3.0,.10)
  {\textcolor{headerblue}{\textbf{04\enspace Pipeline substitution}}\\[2pt]
   42 clusters; topic-block and\\hierarchical intervals.};
\end{tikzpicture}}
\caption{Scope and evidence flow. Top: the five HSMVideoSumm stages, with
visual-quality selection highlighted. Bottom: four complementary tests of the
operation implementing that criterion. All other pipeline stages remain fixed.}
\label{fig:overview}
\end{figure}

\subsection{Data and units of analysis}\label{sec:data}

The external benchmark uses KADID-10k, KonIQ-10k and LIVE In the Wild (CLIVE).
KADID-10k contains 10,125 distorted images derived from 81 references, with 25
distortion types at five levels and Differential Mean Opinion Scores
(DMOS)~\citep{lin2019kadid}. KonIQ-10k contains 10,073 authentically distorted
images, approximately 1.2 million ratings and Mean Opinion Scores
(MOS)~\citep{hosu2020koniq}; CLIVE
contains 1,162 mobile photographs with authentic distortions
~\citep{ghadiyaram2016clive}. No method is retrained or fine-tuned.

The deployment analysis uses the redundancy stage produced by HSMVideoSumm over
19 of its 21 news topics: 42 multi-member clusters, 93 segment slots and 83
distinct segments. The other two topics contain no multi-member cluster. These
clusters are operational inputs, not proof of equal semantic utility. The
original human redundancy annotations cover only 15 clusters and are retained
as a supplementary sensitivity analysis (Tables~S1, S7, S8 and S11); the
42-cluster set is primary because it is the set encountered by the deployed
pipeline.

Still-image scorers receive frames sampled uniformly at 1 frame per second and
are pooled by the arithmetic mean,
\begin{equation}
S(s)=\frac{1}{n_s}\sum_{i=1}^{n_s}\phi(f_i),
\label{eq:meanpool}
\end{equation}
where $\phi$ is the per-frame score. The deployed operation instead applies its
own key-frame selector, quantizes retained descriptors into a vocabulary of
$K=300$ visual words and sums
\begin{equation}
S_{\mathrm{BoVW}}(s)=\sum_i\sum_{v=1}^{K}
\mathrm{tf}(v,f_i)\,\mathrm{idf}(v).
\label{eq:bovw}
\end{equation}
Sampling, scoring and pooling jointly define the operation occupying the
visual-quality stage, so the primary comparison evaluates complete operations
while holding every other stage fixed. A matched-input control additionally gives
SIFT+BoVW and LIQE the same fixed set of at most ten uniformly spaced 1-fps
frames and arithmetic-mean pooling. For SIFT, common-frame descriptors fit a
seeded per-cluster vocabulary with
$K_{\mathrm{eff}}=\min(300,n_{\mathrm{descriptors}})$; candidate segments are
the documents for the deployed $\log_{10}(K_{\mathrm{eff}}/\mathrm{df})$ IDF.
Empty-descriptor frames remain as zeros in the mean; all observed segments
yielded descriptors and no top-score ties occurred. This isolates scorer
identity within the matched-input regime (Tables~S6, S9 and S11).

\subsection{Scorers, references and polarity}\label{sec:methods}

The complete-case comparison contains the deployed SIFT+BoVW score and ten
NR-IQA candidates: BRISQUE, NIQE, DBCNN, HyperIQA, MUSIQ, TOPIQ, CLIP-IQA,
CLIP-IQA$+$, LIQE and QualiCLIP. Keypoint density, Spatial Information, deep
feature norms and the second IDF formulation are diagnostic controls. A
separate ablation replaces SIFT with 19 convolutional extractors while holding
the vocabulary, weighting and aggregation fixed. ARNIQA and TReS extend
coverage on completed axes but are excluded from complete-case screening;
MANIQA was not completed before the experimental freeze (Table~S13).
Released-weight supervision, checkpoint role,
polarity and GPU requirements are summarized in Table~\ref{tab:methods} and
listed in full in Table~S13~\citep{agnolucci2024arniqa,golestaneh2022tres,
yang2022maniqa}.

\begin{table}[!ht]
\centering
\caption{Scorers needed to interpret the main comparison. Sup.: K = released
KonIQ-10k opinion-trained checkpoint; L = LIVE opinion training; OU = opinion-
unaware; M = multi-database opinion mixture; -- = no opinion data. Polarity is
declared before evaluation. D = deployed, C = candidate, R = reference, X =
diagnostic or ablation. Partial-coverage methods remain outside screening.}
\label{tab:methods}
\scriptsize
\setlength{\tabcolsep}{3.3pt}
\renewcommand{\arraystretch}{1.05}
\begin{tabular}{@{}L{2.2cm}L{2.35cm}L{3.65cm}cccc@{}}
\toprule
\rowcolor{headerblue}
\hd{Method} & \hd{Family} & \hd{Measured construct} & \hd{Sup.} & \hd{Pol.} & \hd{GPU} & \hd{Role} \\
\midrule
\rowcolor{catbovw!12}SIFT+BoVW & handcrafted retrieval & TF-IDF-weighted visual-word count & -- & \polup & no & D \\
Keypoint density & handcrafted & detectable keypoints per frame & -- & \polup & no & X \\
\rowcolor{rowlight}DINOv2 & self-supervised features & mean norm of patch tokens & -- & \polup & yes & X \\
Spatial Information & ITU-T complexity & Sobel energy & -- & \polup & no & X \\
\rowcolor{rowlight}BRISQUE & natural-scene statistics & deviation from learned natural statistics & L & \poldown & no & C \\
NIQE & natural-scene statistics & distance from a pristine model & OU & \poldown & no & C \\
\rowcolor{rowlight}DBCNN & CNN regression & bilinear distortion/content features & K & \polup & yes & C \\
HyperIQA & CNN regression & content-adaptive quality regression & K & \polup & yes & C \\
\rowcolor{rowlight}MUSIQ & transformer & multi-scale image quality & K & \polup & yes & C \\
TOPIQ & transformer & top-down quality attention & K & \polup & yes & C \\
\rowcolor{rowlight}CLIP-IQA & vision--language & antonym-prompt quality & OU & \polup & yes & C \\
CLIP-IQA$+$ & vision--language & learned-prompt quality & K & \polup & yes & C \\
\rowcolor{rowlight}QualiCLIP & vision--language & degradation-ranked quality & OU & \polup & yes & C \\
LIQE & vision--language & quality, distortion and scene & K & \polup & yes & C \\
\rowcolor{rowlight}Q-Align & multimodal language model & discrete textual quality levels & M & \polup & yes & R \\
CNN+BoVW (19 extractors) & descriptor ablation & TF-IDF-weighted visual-word count & -- & \polup & yes & X \\
\rowcolor{rowlight}ARNIQA; TReS & partial coverage & distortion manifold; ranking consistency & K & \polup & yes & C \\
\midrule
\multicolumn{7}{@{}L{12.15cm}}{\scriptsize The 19-extractor ablation comprises
ResNet-18/50/101, VGG-16/19, EfficientNet-B0/B4, MobileNet-V3-Large/V2,
ConvNeXt-Tiny/Small, DenseNet-121/201, ResNeXt-50/101,
Wide-ResNet-50/101, Inception-V3 and SqueezeNet-1.1. Released weights,
endpoints and preprocessing are in Table~S13.} \\
\bottomrule
\end{tabular}
\end{table}

Q-Align is an in-domain reference model rather than a candidate
~\citep{wu2024qalign}. The
\texttt{q-future/one-align} checkpoint is run with the \texttt{quality} task,
\texttt{image} input, 4-bit NormalFloat4 (NF4) double quantization and FP16
compute, then pooled by Eq.~\ref{eq:meanpool}. It was among the strongest public NR-IQA models when the
inventory was frozen, is architecturally distinct from the candidates, and is
13.7 times slower than LIQE under the measured setup; it is therefore reserved
as a reference model rather than screened as an on-line candidate.

DOVER is used only as an output-sensitivity reference and does not participate
in screening~\citep{wu2023dover}. The official \texttt{DOVER.pth} checkpoint and
\texttt{dover.yml} \texttt{val-l1080p} views are used (Table~S13). Its raw
technical branch is primary; the uncalibrated technical--aesthetic sum is only a
sensitivity analysis, and neither the official standardized fusion nor DOVER++
is computed. The same technical output is tested on the 800 constructed pairs,
using a common 5-fps MP4 wrapper and paired fragment seed 42. Because DOVER was
developed for user-generated video, its output is not treated as broadcast-domain
human ground truth.

Score and target directions are declared before analysis. Direction-aligned
correlation is $\rho^{*}=s_m s_y\rho$, where $s_m=+1$ for a higher-is-better
method and $-1$ for a lower-is-better method, while $s_y=+1$ when a higher
target means better quality and $-1$ when it means worse quality. All opinion
targets used here---including KADID's released \texttt{dmos} field---are
higher-is-better; degradation level is higher-is-worse. A negative $\rho^{*}$
therefore denotes a genuine directional inversion rather than a scale
convention.

\subsection{Controlled degradation and constructed truth}\label{sec:controlled}

KADID sensitivity is computed without treating derivatives of the same source
as independent. For each distortion type, Spearman correlation between score
and level is computed within each of the 81 references and averaged across
references. Confidence intervals use $10^4$ block-bootstrap resamples of the
reference images with seed 42.

The broadcast experiment supplies a second controlled test. Eight degradations
are applied at five levels: Gaussian blur, white noise, darkening, JPEG
compression, low-bitrate H.264 transcoding, a
$720\mathrm{p}\rightarrow480\mathrm{p}\rightarrow720\mathrm{p}$ round trip,
horizontal motion blur and progressive desaturation. The last uses
$S'_{\ell}=\operatorname{clip}(\alpha_{\ell}S,0,255)$ with
$\alpha_{\ell}\in\{0.75,0.55,0.35,0.18,0.05\}$. H.264 is applied to the frame
sequence; other transformations are frame-based.

Twenty of the 83 broadcast segments are sampled without replacement by a seeded
simple random draw (seed 42); they represent 13 of the 19 operational topics
(Table~S2). Pairing each unmodified segment with every degradation and level gives
$20\times8\times5=800$ binary decisions. The unmodified member is the
construction-defined higher-fidelity signal, not a human perceptual judgment.
Accuracy and paired differences are block-bootstrapped by segment
($10^4$ resamples, seed 42). A second design correlates score with degradation
level within each segment--type pair. Overlay strata, the complete audit of
predictions recorded in a commit whose metadata predates the experiment, and
parameters are reported in Tables~S2 and S12 and the
reproducibility archive. Ties count as neither correct nor incorrect and remain
in the denominator. A sensitivity analysis excludes darkening and desaturation,
whose transformations define signal change more directly than universal
perceptual worsening. These pairs test fidelity response, not preference between
distinct broadcasts.

\subsection{Screening and application inference}\label{sec:screen}

The method-screening rule is operational and was not preregistered. Sensitivity
passes when the within-reference mean is at least .50 and at least 15 of 25
distortion types exceed .50. Generalization passes when the direction-aligned
Spearman Rank Correlation Coefficient (SRCC) is at least .65 on \emph{every}
database absent from the released checkpoint's
opinion training. The in-domain floor is confirmed when normalized agreement
with Q-Align is above chance and the lower bound of its 95\% topic-block
bootstrap interval exceeds zero on the common 42 clusters. Cost is reported but
not thresholded. Table~S10 varies the numerical thresholds and the quantifier
over unseen databases.

For cluster $c$ with $k_c$ candidates, independent uniform choices agree with
probability $1/k_c$. The common-set floor is therefore
$\bar p_e=|C|^{-1}\sum_{c\in C}1/k_c=.470$. Normalized agreement is
$\kappa_{\mathrm{uniform}}=(p_o-\bar p_e)/(1-\bar p_e)$. This is a custom
uniform-choice correction, not Cohen's kappa. A positive point estimate is
reported separately from a confirmed pass so that a wide interval crossing
chance cannot be read as affirmative evidence.

External association uses Spearman correlation against MOS or DMOS. Confidence
intervals (CIs) for KonIQ and CLIVE use the Fisher transform; KADID intervals and paired differences
use reference-block bootstrap. Dependent correlations on the same images use
Steiger's statistic~\citep{steiger1980comparing}. The operational equivalence
margin is $|\Delta\rho^{*}|<.02$, chosen as the largest association difference
treated as practically negligible for this diagnostic rather than as a
literature-wide standard; the paired difference and interval are also reported
so stricter margins can be applied. Within each database, 630 pairwise Steiger
contrasts cover 36 scores (17 methods and 19 ablations); Bonferroni correction
gives $\alpha=.05/630$ and critical $|z|=3.95$. Outside that family, the sole
canonical-score--keypoint equivalence diagnostic uses two one-sided tests at
$\alpha=.05$ against $\pm.02$; KADID also reports its paired block interval.
The full accuracy, consistency and linearity
triad after the monotonic mapping required by ITU-T P.1401 is reported in
Table~S4~\citep{itu_p1401}.

For each operation, the application selects one representative per cluster. If
$y_c$, $l_c$ and $h_c$ are the selected, lowest and highest reference scores in
cluster $c$, define $\bar y=|C|^{-1}\sum_{c\in C}y_c$ and analogously
$\bar l$ and $\bar h$. The reported global-achievable-range estimand is
\begin{equation}
100\,\frac{\bar y-\bar l}{\bar h-\bar l}.
\label{eq:globalrange}
\end{equation}
It is a ratio of means, not an average of independently rescaled clusters.
Raw branch-score differences, mean per-cluster position, selected-maximum rate,
range distributions and exclusions of the smallest 10\% and 25\% of ranges are
reported in Table~S11. Zero-range clusters are undefined for per-cluster
position and are counted explicitly. The primary interval resamples the 19
topics and retains their clusters. A hierarchical sensitivity analysis
additionally resamples clusters within each sampled topic. Both use $10^4$
paired resamples and seed 42. Human compositional preference was not collected;
the estimand is explicitly a model-reference output.

Experiments ran on one NVIDIA RTX~4060 GPU (8~GB). NR-IQA implementations use
the PyIQA toolbox~\citep{chen2022pyiqa}; SIFT uses OpenCV and the vocabulary uses
mini-batch $K$-means with $K=300$ and \texttt{random\_state=42}. Q-Align is
quantized to 4-bit NormalFloat4 with double quantization and FP16 compute to fit
the GPU~\citep{dettmers2023qlora}; the official PyIQA architecture supports
fp16, 4-bit and 8-bit inference, with NF4 used in 4-bit mode. No numerical
equivalence to fp16 is assumed; the reported eligibility set is therefore
conditional on this documented NF4 reference configuration. Frames
are decoded once and cached, and latency is re-measured in one software
environment after ten warm-up and sixty timed iterations. The archive includes
the scoring scripts, derived results, cluster definitions, annotations, seeds
and environment information; redistribution of source news video remains
subject to the broadcasters' terms. The final method inventory and analysis
archive were frozen on 23 August 2026. The historical inference environment was
not preserved as a fully pinned lockfile; the archive therefore reconstructs
the reported statistics from hashed score caches and exposes the inference
settings, but does not claim bit-identical regeneration of every neural score.
